\clearpage
\SourceLeaf{398}
\fancyhead[L]{\small ZUSÄTZE ZUR GEOMETRIE DER STELLUNG.}
\pdfbookmark[1]{Zusätze zur Geometrie der Stellung von Carnot. 1810.}{carnot-zusaetze}
\begin{center}
{\Large ZUSÄTZE.}\par\bigskip
ZUR GEOMETRIE DER STELLUNG VON CARNOT\par\bigskip
ÜBERSETZT VON SCHUMACHER.\quad 1810.\par\bigskip
\rule{.1\textwidth}{.3pt}
\end{center}
\par\medskip\Needspace{5\baselineskip}
\pdfbookmark[2]{I.}{carnot-zusatz-i}
\centerline{\textbf{I.}}\nobreak\smallskip
\{Die folgende analytische Behandlung der bemerkenswerten Punkte eines Dreiecks verdanke ich der Güte von Professor \textsc{Gauss}. \textsc{Schumacher}.\}

Es seien \(A,A',A''\) die drei Eckpunkte eines Dreiecks und ihre jeweiligen Koordinaten
\[
\begin{array}{ll}
x,&y\\x',&y'\\x'',&y''
\end{array}
\]
\par
Die Koordinaten der Punkte \(B,B',B''\), die die Seiten \(A'A'',A''A,AA'\) halbieren, sind offenbar
\[
\begin{array}{ll}
\tfrac12(x'+x''),&\tfrac12(y'+y'')\\
\tfrac12(x''+x),&\tfrac12(y''+y)\\
\tfrac12(x+x'),&\tfrac12(y+y')
\end{array}
\]
\par
Auf den Geraden \(AB,A'B',A''B''\), die bei Bedarf nach vorn oder hinten verlängert werden, trage man von \(A,A',A''\) aus Strecken ab, die zu den jeweiligen Strecken vom Eckpunkt zur gegenüberliegenden Seitenmitte proportional sind und sich zu ihnen wie \(n:1\) verhalten. Trägt man sie nach hinten ab, so ist \(n\) als negativ anzusehen. Die Endpunkte dieser Strecken seien \(C,C',C''\). Ihre Koordinaten lauten
